\begin{abstract}
Large language models achieve remarkable code generation accuracy but struggle to fix their own compiler errors during self-repair, with failure rates exceeding 60\% when given raw error messages. Prior work has focused on \emph{whether} to include compiler feedback, while the question of \emph{how} to format that feedback remains unexplored. We introduce structured error formatting---organizing compiler output with explicit section labels (PROBLEM/LOCATION/CONTEXT/ROOT\_CAUSE)---and demonstrate that this representational alignment significantly improves repair success. Our key insight is that compiler errors fall outside LLM training distributions; reformatting them toward natural language structure enables better model comprehension. Experiments on EvalPlus benchmarks show structured formatting outperforms scrambled formatting with identical content by 14.4 percentage points ($p < 0.001$), providing causal evidence that \emph{how} information is organized matters, not just \emph{what} information is present. Additionally, intermediate-specificity fix hints achieve optimal repair (60.2\%), outperforming both no hints and exact fixes, consistent with scaffolding theory. Our findings suggest that interface design between tools and models deserves the same attention as model architecture itself.
\end{abstract}
